\documentclass[11pt]{article}

% ---------------------------------------------------------------------------
% AfriMentor AI — arXiv preprint scaffold (G1.6)
% Not tied to a specific venue's style file (no acl.sty/neurips.sty assumed
% available) — a clean, standard article-class preprint, which is what
% arXiv itself expects as a baseline. Swap in a venue style later if the
% team decides to target a specific workshop/conference.
% ---------------------------------------------------------------------------

\usepackage[margin=1in]{geometry}
\usepackage{amsmath,amssymb}
\usepackage[colorlinks=true,linkcolor=blue,citecolor=blue,urlcolor=blue]{hyperref}
\usepackage{natbib}
\usepackage{graphicx}
\usepackage{authblk}
\usepackage{booktabs}
\usepackage{times}

\title{AfriMentor AI: A Personality-Aligned, RLHF-Tuned Language Model Mentor for Low-Income African Youth Entrepreneurs}

% Placeholder author block — confirm exact order, affiliations, and
% corresponding-author designation with the team before submission.
\author[1]{Marie Grace Kagaju}
\author[1]{Olusegun Banji Kolewole}
\author[1]{Daniel Kusi Boateng}
\author[1]{Olisa Martin Chukwuebuka}
\affil[1]{AfriMentor AI Project}

\date{\today}

\begin{document}
\maketitle

\begin{abstract}
% TODO (team): write after Sections 3-4 (Method/System, Evaluation) exist —
% an abstract written before the results section is speculative by
% definition. Placeholder left here so the document structure is complete
% and compiles; do not submit with this section empty.
This is a placeholder abstract. AfriMentor AI is a smartphone-first,
voice-note-driven mentorship platform that pairs African micro-entrepreneurs
with an AI mentor persona whose reasoning style is constrained to match a
documented archetype of real African achievers, trained/aligned via
personality-conditioned RLHF, and grounded in an Africa-authored knowledge
corpus. \textbf{[Expand once Sections 3-4 exist; state the two research
questions and headline findings explicitly.]}
\end{abstract}

\section{Introduction}
\label{sec:introduction}

Financial mentorship is one of the strongest predictors of small-business survival, yet it remains one of the least accessible resources for young entrepreneurs across much of Africa. Formal mentorship programs are concentrated in a handful of urban hubs, typically require an existing professional network to access, and are rarely available in the languages, financial contexts, or business categories most relevant to informal-sector entrepreneurs -- market traders, smallholder agribusiness owners, and micro-manufacturers who make up a substantial share of the continent's economic activity. At the same time, smartphone penetration across the region continues to rise even where broadband infrastructure and formal financial education do not, creating an opening for AI-mediated mentorship to reach entrepreneurs that traditional programs cannot.

Large language models (LLMs) are an increasingly plausible vehicle for this kind of mentorship at scale, but a generic conversational assistant is a poor substitute for a mentor. Mentorship is not simply information delivery; it is guidance filtered through a specific, credible point of view -- the way a mentor who has personally built something from constrained circumstances reasons about risk, setbacks, and resource allocation is qualitatively different from, and more persuasive than, an undifferentiated helpful-assistant voice. This motivates the central design question behind AfriMentor AI: \emph{can a language model be constrained to consistently reason in the style of a specific, real archetype of achievement, in a way that is both measurably faithful to that archetype and genuinely useful to the person receiving guidance from it?}

We approach this question through two complementary research threads, which this paper reports on jointly. The first (Proposal 1) is methodological: we investigate whether reinforcement learning from human feedback (RLHF), conditioned on structured personality-trait profiles derived from documented African achievers, can produce a language model whose reasoning style remains consistent across a full mentorship conversation -- not just on a single prompt -- and whether such a model's outputs are distinguishable, in blind evaluation, from a generic same-capability baseline. This connects to a growing body of work on RLHF as a general alignment technique \citep{bai2022helpful,chaudhari2024rlhf} and on measuring whether personality-like traits in LLM outputs are behaviorally consistent rather than merely a consistent \emph{self-report} that fails to predict actual behavior \citep{lee2024trait,han2025illusion,jiang2023personallm}, but applies both specifically to the constraint of fidelity to a \emph{real, individual} archetype rather than a synthetic or aggregate persona -- closer in spirit to recent multi-turn reinforcement learning approaches that measurably reduce persona drift across a full conversation \citep{abdulhai2025personas} than to single-turn role-play.

The second thread (Proposal 2) is empirical and user-centered: what do African micro-entrepreneurs, specifically those with limited formal education and constrained access to data and devices, actually want from an AI mentorship product, and does a personality-constrained persona change engagement or perceived usefulness relative to a generic assistant? This question sits within a broader literature on LLM-based tutoring and educational support systems, much of which has been developed and evaluated through conversation-based tutoring architectures with explicit student modeling and retrieval-augmented content \citep{park2024tutoring,liu2025lpitutor}, and through systematic reviews highlighting both the promise and the open ethical and accessibility questions of LLM-based personalized learning \citep{sharma2025review}. We are specifically interested in whether those findings transfer to an informal-economy, low-bandwidth, voice-first context, or whether they need to be substantially rethought.

The system built to support this research, AfriMentor AI, is a smartphone-first mentorship platform with three defining constraints, chosen specifically to make both research threads testable rather than purely aspirational: (1) a knowledge base restricted to African-authored financial and entrepreneurship content, so that guidance is contextually grounded rather than generically Western; (2) a voice-note-first interaction model, chosen to match the primary user's actual device and connectivity constraints rather than assuming typed, always-online interaction; and (3) an explicit, inspectable mapping from a small set of achiever archetypes to persona-conditioning parameters, so that persona fidelity is something we can actually measure rather than only claim.

The remainder of this paper is organized as follows. Section~\ref{sec:related-work} situates this work relative to prior research on RLHF alignment, personality modeling in LLMs, generative agents, and LLM-based tutoring. Section~\ref{sec:method} (forthcoming) describes the persona-conditioning and RLHF methodology in detail. Section~\ref{sec:evaluation} (forthcoming) reports the persona-fidelity and user-study evaluation. We close with a discussion of limitations specific to deploying this kind of system in low-resource, informal-economy contexts.

\section{Related Work}
\label{sec:related-work}

We situate AfriMentor AI relative to four bodies of work drawn directly from the project's two research proposals: reinforcement learning from human feedback as a general alignment technique; measuring whether personality-like traits in LLMs are behaviorally real or only a consistent self-report; mechanisms for controlling and steering persona expression once it is measured; and LLM-based tutoring and personalized learning. \textbf{A correction, for the record:} an earlier draft of this section cited several papers found via independent web search that matched an author surname and publication year given in project planning materials, but turned out to be the wrong paper -- a different, unrelated work by a different author with the same last name. That draft has been discarded. Every citation below is instead drawn directly from the reference lists in the team's own Proposal 1 and Proposal 2 documents, which is the correct authoritative source.

\subsection{Reinforcement learning from human feedback as an alignment technique}

RLHF is now the standard technique for steering language model behavior toward what human raters judge helpful and appropriate. \citet{bai2022helpful} established the canonical version of this pipeline -- training a reward model from human preference comparisons and using it to reinforcement-learn a policy -- and showed it improves general helpfulness and harmlessness without degrading other capabilities. \citet{chaudhari2024rlhf} take stock of the technique's maturity since then, offering a critical survey of RLHF's mechanics, failure modes, and open methodological questions across the growing body of work that has built on it. Two extensions matter specifically for a personality-constrained system like AfriMentor AI's: \citet{dai2023saferlhf} decouple the helpfulness and harmlessness objectives into a constrained optimization rather than a single blended reward, which is directly relevant to a mentor persona that must stay in-character while also never crossing into unsafe financial advice; and \citet{zhang2025mmrlhf} extend the RLHF framework to multimodal alignment, relevant background given AfriMentor AI's voice-note-first interaction model.

Despite this maturity, none of these methods were designed with an explicit, human-interpretable personality profile as the alignment target -- the reward signal in each case is a general notion of quality or safety, not fidelity to a specific character. That gap is exactly what Proposal 1 addresses.

\subsection{Measuring personality in language models: self-report versus behavior}

A separate line of work asks whether LLMs exhibit personality-like traits at all, and -- more contentiously -- whether traits measured through self-report actually predict how the model behaves in realistic tasks. \citet{lee2024trait} introduce a psychometrically grounded benchmark (TRAIT) built on validated human personality instruments, and use it to show that LLMs express personality profiles that are both distinct between models and stable within a model, shaped substantially by the model's training and alignment data. \citet{jiang2023personallm} study a related but narrower question specifically for role-play: whether an LLM assigned a persona can express that persona's intended traits in its generated text, and find that human raters can recognize the intended trait from LLM output at rates well above chance. \citet{li2024big5chat} approach personality from the training side rather than the measurement side, showing that fine-tuning on dialogue data grounded in real human personality profiles can shift a model's expressed personality in a targeted, controllable direction.

\citet{han2025illusion} complicate all three of these results in a way that is central to how AfriMentor AI's own evaluation must be designed. Their finding is that instructional alignment (of the kind RLHF and instruction-tuning both perform) makes a model's \emph{self-reported} personality more internally consistent -- but that consistency does not reliably transfer to the model's actual behavior in realistic, non-self-report tasks. In other words, a model can pass a personality questionnaire convincingly while still not behaving like the persona it claims in an actual conversation. This is a direct warning against evaluating AfriMentor AI's persona fidelity solely by asking the model to self-report its own traits (e.g., "are you being consistent with the Market Queen archetype right now?") -- the self-report itself is exactly the measurement this literature shows can be misleading, which is why Proposal 1's evaluation plan pairs trait-level self-report scoring with separate behavioral consistency metrics rather than relying on either alone.

\subsection{Controlling and steering persona expression}

A third thread asks not just whether personality can be measured, but how it can be deliberately controlled once measured, at several different levels of the model. \citet{chen2025personavectors} and \citet{frising2025linear} both work at the level of internal model representations, identifying directions in a model's activation space that correspond to specific personality dimensions and showing these can be used to monitor or directly steer trait expression without retraining the model. \citet{kose2025hallucination} apply a similar persona-vector approach with a safety framing, examining whether steering a small model's persona representation can also reduce its tendency to hallucinate. At the level of training rather than inference-time steering, \citet{ji2025pcl} propose a contrastive learning objective specifically designed to keep a role-playing model's outputs consistent with its assigned persona's characteristics across a conversation, rather than drifting toward generic responses. \citet{frisch2024agents} study a related consistency question empirically, measuring how stable an LLM agent's expressed personality and linguistic style remain when it interacts with other LLM agents over an extended exchange, rather than a single isolated prompt.

\citet{abdulhai2025personas} is the closest single precedent to AfriMentor AI's own methodological approach in Proposal 1: they define automatic, reference-based consistency metrics for persona fidelity across a multi-turn conversation (comparing what a persona says early in a conversation against what it says later, rather than only checking each response in isolation), and show that using these metrics as a training signal in a multi-turn reinforcement learning loop substantially reduces persona drift compared to prompting alone. AfriMentor AI's Proposal 1 adopts a structurally similar approach -- multi-turn, reference-grounded consistency metrics used as an RLHF signal -- applied to the specific case of fidelity to one documented individual archetype rather than an open-ended role-play character.

Finally, \citet{rahwan2019machine} frame the broader stakes of this entire line of work: as autonomous and semi-autonomous AI systems increasingly exhibit consistent, characterizable behavioral patterns, understanding and predicting those patterns becomes a legitimate empirical science in its own right, not merely an engineering afterthought. AfriMentor AI's persona-fidelity evaluation is, in that framing, a domain-specific instance of exactly this broader "machine behaviour" research agenda.

\subsection{LLM-based tutoring and personalized learning}

The second research thread (Proposal 2) sits within a distinct but related literature on LLMs as tutors and personalized learning systems. \citet{park2024tutoring} present a conversation-based tutoring system that incorporates explicit diagnostic modeling of the learner -- tracking what the student already knows, not just what they most recently asked -- and find this diagnostic layer meaningfully improves how personalized the system's responses feel. \citet{liu2025lpitutor} build a comparably structured system (LPITutor) that combines retrieval-augmented generation with prompt engineering to balance response accuracy against clarity, adapting explanation complexity to the difficulty level implied by the learner's own question. \citet{sharma2025review} step back from any single system to synthesize the field, systematically reviewing LLM-based personalized learning research and concluding that while these systems show real promise for engagement, real-time monitoring, and emotional support, the literature still has open, largely unresolved questions around ethics, bias, privacy, and whether LLMs can authentically reproduce pedagogical strategies -- such as Socratic questioning -- that depend on more than surface-level dialogue fluency. \citet{grassucci2025strategic} examine a related but distinct capability question: whether LLMs used in education can support genuine strategic thinking in learners, rather than only fluent explanation.

Both \citet{park2024tutoring}'s diagnostic student modeling and \citet{liu2025lpitutor}'s retrieval-augmented, prompt-engineered architecture are direct methodological precedents for AfriMentor AI's own design -- Proposal 2's system likewise tracks a simple user model (current knowledge, stated goals, and past actions) and layers a curated retrieval corpus underneath persona-conditioned prompting. The open questions \citet{sharma2025review} raise around ethics, privacy, and authenticity of pedagogical strategy are directly inherited by AfriMentor AI's context: privacy is a first-class concern given the platform's research data collection from a vulnerable, low-income user population, and authenticity of reasoning style -- not just tone -- is precisely what Proposal 1's persona-fidelity metrics attempt to measure rather than assume.

Notably, this tutoring and personalized-learning literature is concentrated in formal or semi-formal educational settings with a defined curriculum and, typically, reliable connectivity and typed interaction. None of it is situated in the specific constraint set that most directly shapes AfriMentor AI's design: an informal-economy user base, voice-note-first interaction chosen to match actual device and data constraints rather than assumed always-online typing, and financial (not academic) guidance as the subject matter. AfriMentor AI sits at the intersection the two source literatures do not individually cover -- personality-constrained, RLHF-aligned mentorship, evaluated specifically among low-income African entrepreneurs under real bandwidth and device constraints -- which is the gap this paper's combined contribution is intended to address.


% Not yet drafted — stubbed so \ref/\cite structure and the build stay
% complete. See contract with the team on the 1-paper vs 2-paper decision
% (docs/paper-scope-decision.md) before writing further, since it changes
% how much lives under each of these.
\section{Method}
\label{sec:method}
\textit{[Not yet drafted. Depends on Epic H (trait modeling) and Epic C
(persona-constrained personas) reaching enough maturity to describe
concretely — both are mid-sprint as of this draft.]}

\section{Evaluation \& Results}
\label{sec:evaluation}
\begin{table}[ht]
\centering
\caption{Comparative evaluation of the four alignment conditions across the shared metric suite. Scores are means over the held-out test set (\texttt{sft\_test.jsonl}, $N=5$). Composite is the weighted sum of the five rubric dimensions. All four conditions are live measurements (\texttt{source}: \texttt{live\_groq} for C1, \texttt{live\_hf\_adapter} for C2--C4) from the frozen evaluation snapshot \texttt{eval-freeze-v1} — see Appendix~\ref{sec:appendix-reproducibility}.}
\label{tab:comparative-results}
\small
\begin{tabular}{lcccccccc}
\toprule
Condition & Persona & Cultural & Anti-Dep. & Financial & Urgency & ROUGE-L & BERTScore & Composite \\
& Adherence & Fluency & & Accuracy & & & F1 & \\
\midrule
C1: Baseline Prompting & 0.780 & 0.560 & 0.710 & 0.640 & 0.490 & 0.138 & 0.821 & 0.651 \\
C2: Supervised Fine-Tuning & 0.890 & 0.610 & 0.710 & 0.800 & 0.710 & 0.134 & 0.849 & 0.753 \\
C3: Contrastive Learning (DPO) & 0.820 & 0.650 & 0.660 & 0.660 & 0.650 & 0.134 & 0.849 & 0.697 \\
C4: RLHF / Preference Opt. & 0.700 & 0.590 & 0.650 & 0.640 & 0.610 & 0.134 & 0.850 & 0.643 \\
\bottomrule
\end{tabular}
\end{table}


\section{Discussion \& Limitations}
\label{sec:discussion}
\textit{[Not yet drafted.]}

\section{Conclusion}
\label{sec:conclusion}
\textit{[Not yet drafted.]}

\bibliographystyle{plainnat}
\bibliography{refs}

\end{document}
